\documentclass[a4paper,12pt]{article}

\usepackage[T2A]{fontenc}
\usepackage[utf8]{inputenc}
\usepackage[russian]{babel}
\usepackage{amsmath,amssymb,mathtools}
\usepackage{PTSans}
\renewcommand{\familydefault}{\sfdefault}
\usepackage{icomma}
\usepackage[a4paper,margin=19mm,headheight=25pt]{geometry}
\usepackage{microtype}
\usepackage{tikz}
\usetikzlibrary{calc,arrows.meta,positioning,shapes.geometric}
\usepackage{tabularx,booktabs}
\usepackage{enumitem}
\usepackage{fancyhdr}
\usepackage{setspace}
\usepackage{xcolor}
\usepackage{array}

\definecolor{accent}{HTML}{2B6F6D}
\definecolor{darkaccent}{HTML}{184B49}
\definecolor{textgray}{HTML}{2F3437}
\definecolor{softgray}{HTML}{F3F5F4}

\geometry{left=19mm,right=19mm,top=19mm,bottom=20mm}
\onehalfspacing
\setlength{\parindent}{0pt}
\setlength{\parskip}{0.45em}
\setlist[itemize]{leftmargin=1.55em,itemsep=0.22em,topsep=0.2em}
\setlist[enumerate]{leftmargin=1.8em,itemsep=0.55em,topsep=0.35em}

\pagestyle{fancy}
\fancyhf{}
\lhead{\small\color{textgray}Финансово-экономические задачи}
\rhead{\small\color{textgray}Тимофей Васильченко}
\cfoot{\small\color{textgray}\thepage}
\renewcommand{\headrulewidth}{0.3pt}
\renewcommand{\headrule}{\hbox to\headwidth{\color{accent}\leaders\hrule height \headrulewidth\hfill}}

\newcommand{\sectiontitle}[1]{%
  \vspace{0.4em}
  {\Large\bfseries\color{darkaccent}#1}\par
  \vspace{0.15em}
  {\color{accent}\rule{\linewidth}{0.6pt}}\par
  \vspace{0.25em}
}

\newcommand{\checkitem}{\item[$\square$]}
\newcommand{\rub}{\text{ руб.}}
\newcommand{\thrub}{\text{ тыс. руб.}}
\newcommand{\mlnrub}{\text{ млн руб.}}

\begin{document}

\begin{titlepage}
\thispagestyle{empty}
\centering

\vspace*{18mm}

{\Large\color{accent}\bfseries Чек-лист\par}
\vspace{5mm}
{\Huge\bfseries\color{darkaccent}Финансово-экономические задачи\par}
\vspace{2mm}
{\LARGE\bfseries\color{textgray}Вклады и кредиты\par}

\vspace{14mm}

{\large\color{textgray}
Профильная математика, ЕГЭ\par}

\vspace{18mm}

\begin{tabular}{>{\bfseries}r l}
Ученик: & Тимофей Васильченко\\[2mm]
Преподаватель: & Лёвин Артём Александрович\\[2mm]
Дата: & 27 сентября 2026 г.
\end{tabular}

\vfill

{\large\color{darkaccent}\bfseries
Лёвин Артём Александрович эксклюзивно для Тимофея Васильченко
\par}

\vspace{10mm}

\begin{tikzpicture}[x=1cm,y=1cm]
  \draw[accent, line width=1.2pt]
    (-7.2,0)
    .. controls (-4.9,1.2) and (-2.8,-0.8) .. (-0.7,0.25)
    .. controls (1.6,1.3) and (4.2,-0.9) .. (7.2,0.15);
  \fill[accent] (-7.2,0) circle (1.4pt);
  \fill[accent] (7.2,0.15) circle (1.4pt);
\end{tikzpicture}

\vspace*{5mm}
\end{titlepage}

\sectiontitle{1. Главная идея}

Финансовая задача становится существенно проще, если Вы рассматриваете деньги как величину,
которая \textbf{последовательно меняется во времени}. Для каждого года важно фиксировать
не только процентную ставку, но и \textbf{порядок событий}: когда начисляются проценты,
когда совершается пополнение, снятие или платёж.

\begin{itemize}
\checkitem Переводите процентную ставку в десятичную дробь:
\[
p=\frac{r}{100}.
\]

\checkitem При увеличении суммы на \(r\%\) удобно использовать множитель
\[
1+p=1+\frac{r}{100}.
\]

\checkitem Если сумма \(S\) увеличивается на \(r\%\), то
\[
S_{\text{нов}}=S(1+p).
\]

\checkitem Если величина неизвестна, вводите переменную именно для того, что требуется найти
или что повторяется из года в год: \(x\), \(p\), \(S\).

\checkitem Всегда отделяйте \textbf{тело вклада/кредита} от процентов, платежей и дополнительных взносов.
\end{itemize}

\sectiontitle{2. Вклады: таблица и порядок операций}

Для вкладов полезно фиксировать не абстрактные «шаги», а \textbf{события в их реальном порядке}:
\[
\begin{aligned}
S_{\text{нач}}
&\to \text{операция до процентов}
\to \text{база начисления},\\
&\to \text{проценты}
\to \text{операция после процентов}
\to S_{\text{кон}}.
\end{aligned}
\]
Не каждый столбец обязательно заполняется: если операции в соответствующий момент нет, просто переходите к следующему шагу.

Главный вопрос каждого года: \textbf{какая сумма реально лежит на вкладе в момент начисления процентов?}

\begin{itemize}
\checkitem Если пополнение сделано \textbf{в начале года до начисления процентов}, сначала прибавьте его к сумме на счёте, а затем начисляйте процент.

\checkitem Если пополнение сделано \textbf{после начисления процентов}, оно увеличивает итог текущего года, но проценты на эту добавку начнут начисляться только в следующем периоде.

\checkitem Фраза «в конце первого года после начисления процентов внесли \(A\)» эквивалентна фразе
«в начале второго года до начисления процентов внесли \(A\)».

\checkitem При снятии денег действует та же логика: сначала определите момент снятия, затем установите базу, от которой начисляется процент.
\end{itemize}

\textbf{Пример.}
На вклад внесли \(120\mlnrub\). В конце каждого года сумма увеличивается на \(15\%\).
В начале второго и третьего годов вносят ещё по \(10\mlnrub\).

\[
120 \xrightarrow{\,+15\%\,} 138,
\qquad
138+10=148,
\]
\[
148\cdot1{,}15=170{,}2,
\qquad
170{,}2+10=180{,}2,
\]
\[
180{,}2\cdot1{,}15=207{,}23.
\]

К концу третьего года на вкладе будет
\[
\boxed{207{,}23\mlnrub}.
\]

\sectiontitle{3. Неизвестная ставка и меняющиеся проценты}

Если ставка одна и та же каждый год и дополнительных операций нет, можно сразу использовать
степень:
\[
S_n=S_0(1+p)^n.
\]

Если ставка неизвестна, удобно обозначить
\[
x=1+p.
\]
Тогда
\[
S_n=S_0x^n.
\]

\textbf{Пример.}
\(500\thrub\) превратились за три года в \(665{,}5\thrub\):
\[
500x^3=665{,}5,
\qquad
x^3=1{,}331,
\qquad
x=1{,}1.
\]
Следовательно,
\[
p=0{,}1,\qquad r=10\%.
\]

Если процентная ставка меняется по годам, множители также меняются:
\[
S_3=S_0(1+p_1)(1+p_2)(1+p_3).
\]
Например, для ставок \(8\%\), \(10\%\), \(12\%\):
\(S_3=S_0\cdot1{,}08\cdot1{,}10\cdot1{,}12\).

\sectiontitle{4. Кредиты: универсальная таблица}

Для кредита полезно держать одну и ту же структуру:

\begin{center}
\begin{tabularx}{\linewidth}{>{\bfseries}l X}
\toprule
Этап & Что означает \\
\midrule
Начальный долг & Сколько Вы должны банку на начало года \\
Начисленные проценты & \(p\) от начального долга текущего года \\
Общая задолженность & Начальный долг \(+\) начисленные проценты \\
Платёж & Сколько вносится банку в этом году \\
Остаток долга & Общая задолженность \(-\) платёж \\
\bottomrule
\end{tabularx}
\end{center}

Для стандартного порядка «проценты \(\to\) платёж»:
\[
D_{\text{после}}=D_{\text{до}}(1+p)-P,
\]
где \(D_{\text{до}}\) --- долг на начало периода, \(P\) --- платёж.
При другом порядке событий сначала перестройте строку таблицы по условию.

\begin{itemize}
\checkitem Остаток одного года становится начальным долгом следующего.

\checkitem Если к концу срока кредит полностью погашен, последний остаток равен нулю.

\checkitem Именно условие \(D_{\text{после}}=0\) часто даёт уравнение для неизвестного платежа.

\checkitem Схема погашения определяется текстом задачи. Нельзя автоматически считать платежи одинаковыми.
\end{itemize}

\sectiontitle{5. Три важные схемы погашения кредита}

\textbf{Схема 1. Платёж равен начисленным процентам.}

Если после платежа долг должен остаться прежним, то платёж компенсирует только проценты:
\[
P=pD.
\]
Тело кредита при этом не уменьшается.

\medskip

\textbf{Схема 2. Платежи одинаковы.}

Если по условию несколько платежей равны, обозначайте каждый из них одной переменной:
\[
P_1=P_2=\dots=x.
\]
После заполнения таблицы используйте условие полного погашения кредита.

\medskip

\textbf{Схема 3. Долг уменьшается равными долями.}

Если кредит \(S\) взят на \(n\) лет и тело долга ежегодно уменьшается на одинаковую величину,
то ежегодно погашается
\[
\frac{S}{n}
\]
\textbf{именно от исходной суммы кредита}.

Платёж каждого года состоит из двух частей:
\[
P_k=\frac{S}{n}+pD_k,
\]
где \(D_k\) --- долг на начало соответствующего года.

\textbf{Пример.}
Кредит \(120\thrub\) взят на 3 года под \(10\%\), долг уменьшается равными долями.
Ежегодное погашение тела:
\[
\frac{120}{3}=40\thrub.
\]
Поэтому платежи:
\[
40+12=52,\qquad
40+8=48,\qquad
40+4=44.
\]
Остаток долга проходит цепочку
\[
120\to80\to40\to0.
\]

\sectiontitle{6. Частые ошибки}

\begin{itemize}
\checkitem \textbf{Проценты считают не от той суммы.}
Сначала установите, сколько денег реально находится на вкладе или сколько долга реально осталось
в момент начисления процентов.

\checkitem \textbf{Путают проценты и коэффициент роста.}
\(10\%=0{,}1\), а множитель увеличения равен \(1{,}1\).

\checkitem \textbf{Путают равные платежи и равномерное уменьшение долга.}
При равных долях тела кредита сами платежи обычно различаются, потому что проценты начисляются на уменьшающийся долг.

\checkitem \textbf{Забывают последний нулевой остаток.}
Если кредит полностью погашен, финальная строка таблицы должна приводить к уравнению
\[
D_{\text{после}}=0.
\]

\checkitem \textbf{Рано округляют.}
В длинных расчётах сохраняйте точные десятичные значения до финального ответа.

\checkitem \textbf{Сразу пишут громоздкую формулу.}
Сначала заполните таблицу по годам. После этого уравнение обычно возникает естественно.
\end{itemize}



\sectiontitle{7. Тренировочный блок --- Путь А}

\textbf{10 задач.} Решайте через таблицу или последовательность состояний.
В тренировочном блоке решений нет.

\begin{enumerate}
\item На вклад внесли \(200\thrub\). В конце каждого года банк увеличивает сумму вклада на \(12\%\).
Дополнительных операций нет. Какая сумма будет на вкладе через два года?

\item На вклад внесли \(120\thrub\). В конце каждого года сумма увеличивается на \(10\%\).
В начале второго и третьего годов со вклада снимают по \(5\thrub\).
Сколько будет на вкладе к концу третьего года?

\item На вклад внесли \(100\thrub\). В конце каждого года сумма увеличивается на \(20\%\).
В начале второго и третьего годов вклад пополняют на \(10\thrub\).
Найдите сумму на вкладе к концу третьего года.

\item Вклад \(400\thrub\) размещён на три года без пополнений и снятий.
Процентная ставка каждый год одна и та же.
К концу третьего года на счёте оказалось \(532{,}4\thrub\).
Найдите годовую процентную ставку.

\item На вклад внесли \(1000\thrub\).
В конце первого года начислили \(8\%\), в конце второго --- \(10\%\),
в конце третьего --- \(12\%\).
Пополнений и снятий не было.
Найдите итоговую сумму.

\end{enumerate}

\clearpage
\textbf{Продолжение тренировочного блока.}
\begin{enumerate}[resume]

\item Кредит \(240\thrub\) взят на три года под \(10\%\) годовых.
После каждого ежегодного платежа тело долга уменьшается на одну и ту же величину,
а к концу третьего года кредит полностью погашен.
Найдите сумму всех трёх платежей.

\item Кредит \(300\thrub\) взят на три года под \(20\%\) годовых.
В конце первого и второго годов платёж равен только начисленным процентам,
поэтому тело долга остаётся неизменным.
В конце третьего года кредит погашается полностью.
Найдите сумму всех платежей.

\item Кредит \(500\thrub\) взят на три года под \(10\%\) годовых.
Проценты начисляются в конце каждого года перед платежом.
Платежи в конце первого и второго годов одинаковы и равны \(x\),
платёж в конце третьего года равен \(203{,}5\thrub\).
После третьего платежа долг равен нулю.
Найдите \(x\).

\item Кредит \(600\thrub\) взят на четыре года под \(15\%\) годовых.
После каждого ежегодного платежа тело кредита уменьшается на одну и ту же величину,
а к концу четвёртого года долг равен нулю.
Найдите сумму всех четырёх платежей.

\item Кредит \(360\thrub\) взят на четыре года под \(25\%\) годовых.
В конце первого и второго годов заёмщик выплачивает только начисленные проценты,
поэтому тело кредита остаётся равным \(360\thrub\).
Платежи в конце третьего и четвёртого годов одинаковы и равны \(x\).
После четвёртого платежа долг равен нулю.
Найдите \(x\).
\end{enumerate}

\clearpage

\sectiontitle{8. Ответы}

\renewcommand{\arraystretch}{1.35}
\begin{center}
\begin{tabular}{>{\bfseries}c l}
\toprule
№ & Ответ \\
\midrule
1 & \(250{,}88\thrub\) \\
2 & \(148{,}17\thrub\) \\
3 & \(199{,}2\thrub\) \\
4 & \(10\%\) \\
5 & \(1330{,}56\thrub\) \\
6 & \(288\thrub\) \\
7 & \(480\thrub\) \\
8 & \(200\thrub\) \\
9 & \(825\thrub\) \\
10 & \(250\thrub\) \\
\bottomrule
\end{tabular}
\end{center}

\vfill

{\small\color{textgray}
Ключевой ориентир самопроверки: в каждом году сначала определите сумму,
от которой начисляются проценты, и только после этого выполняйте следующую операцию.
}

\end{document}
